\chapter{What a Language Presentation Is, and Why It Comes First}
\label{ch:gslt-intro}
% ===================================================================

The rest of this book builds a mind out of computations, and then asks
what such a mind can know, what physics it finds itself inside, and how
anything of the kind could have got started on a planet.
Every one of those arguments spends the same coin.
This part mints it.

The coin is a way of writing down a language --- not a program, not a
machine, but a \emph{language} --- that is regular enough to be handed to
a machine rather than to a human implementer.
A grammar, a set of equations saying which terms are the same, and a set
of rewrite rules saying how one term becomes another.
That triple is a \emph{graph-structured lambda theory}, and the claim of
this part is that once you have one, an astonishing amount comes free:
a notion of interaction, a meter that charges for it, a log that
remembers it, a logic that describes it, and a type system that checks
it.

\section{Why this order}

An earlier arrangement of this material put the physics first.
It began with a rewrite theory, hung weights and costs on the rewrite
steps, recovered least action and conservation, and only much later
introduced anything that could be called an agent --- at which point the
agent arrived as a placeholder, a thing that has theories and can afford
some number of experiments.

That order encodes a claim we no longer believe: that the world comes
first and the knower is a late arrival in it.
The difficulty is not philosophical delicacy.
It is that under the earlier order the physics is presented as though it
were the only one, and the learner as though it had been dropped into
it.
Both halves of that picture are wrong in the same way.
A GSLT does not come with a physics; it comes with a \emph{space} of
possible physics, and which one you get depends on what further axioms
you are willing to impose.
A learner does not arrive in a world; it is made of the same material as
the world and pays the same prices.
Part~\ref{part:physics} makes the first point, and it can only make it
after Part~\ref{part:ecology} has made the second.

So the order here is: machinery, then mind, then the physics the mind
finds itself inside, then what the mind can do that a Turing machine
cannot, then how such a thing could have arisen.
This part is the machinery.
It is deliberately the driest part of the book, and a reader who wants
the argument rather than the apparatus may read
Chapters~\ref{ch:gslt}, \ref{ch:igslt} and~\ref{ch:cigslt} for the three
definitions, skim Chapter~\ref{ch:costmonad} for the shape of the cost
construction, read Chapter~\ref{ch:carho} for what a metered language
actually looks like, and go on to Part~\ref{part:ecology}, returning here
when a construction is invoked that they want to see built.

\section{The three definitions, in advance}

The development proceeds by two strengthenings of a very thin starting
notion, and it is worth saying now what each buys, because the discipline
of adding structure only when a construction needs it makes the sequence
harder to see from inside.

\begin{description}[leftmargin=1.6em,style=nextline]
\item[A GSLT] is any rewrite theory whatever: a signature, equations,
  rewrites.
  Nothing in it says the theory is about computation, or that anything
  interacts with anything.
  This is enough for the category of Chapter~\ref{ch:catgslt} and for the
  type generation of Chapter~\ref{ch:hypercube}.
\item[An interactive GSLT] names the \emph{site} at which two things meet
  --- application in the $\lambda$-calculus, parallel composition in rho ---
  and distinguishes the base rule that fires there from the context rules
  that say when it may.
  This is what lets us speak of an agent and its environment without
  smuggling in a homunculus: the boundary is drawn by the rules, at
  runtime, and either side of an interaction may be the program.
\item[A continued interactive GSLT] presents that site as a
  \emph{cut} which can be metered: the interaction is factored so that
  there is a well-defined place to put a charge, and continuations become
  metered thunks.
  This is what makes Chapter~\ref{ch:costmonad} possible, and with it
  everything in this book that depends on a learner having a budget ---
  which is to say everything in this book.
\end{description}

The examples matter as much as the definitions, including the ones that
fail.
A Turing machine is a GSLT and is not interactive; a $\lambda$-calculus is
interactive with an asymmetric site; rho is interactive with a symmetric
one.
Chapter~\ref{ch:igslt} is largely an argument about which is which and
why the distinction is not a technicality.

\section{Running example: the rho calculus}
\label{sec:rho-running}

Throughout we use the rho calculus~\cite{meredith2005rho} as the running
example, because it exhibits in compact form nearly every feature at
issue: reflection, so that processes and names are mutually definable and
a term can carry a quotation of another term; parallel composition, so
that the interaction site is symmetric; and a single synchronization rule
that will serve as the prototype for causal structure, for replication,
and for the act of one computation eating another.

The grammar is
\begin{align}
  P, Q &\;::=\; 0 \;\mid\; \mathtt{for}(y \leftarrow x)\,P \;\mid\;
               x!(Q) \;\mid\; P \mid Q \;\mid\; {*}x \\
  x, y &\;::=\; @P
\end{align}
where $y$ in $\mathtt{for}(y \leftarrow x)\,P$ is bound.
Structural equivalence $\equiv$ is the smallest equivalence relation
containing $\alpha$-equivalence and making $(P,\,|\,,0)$ a commutative
monoid.
The single reduction rule is
\[
  \mathtt{for}(y \leftarrow x)\,P \;\mid\; x!(Q)
  \;\rewrite\;
  P\{@Q / y\}.
\]

Three features of this calculus do disproportionate work later.
Reflection means the quote of a process is a name, so a term can hold
code it has not run --- which is what makes the genetics of
Chapter~\ref{ch:com} possible, since crossover is an operation on quoted
code.
Names are the only locations there are, so a namespace is a region, and
a logic of namespaces is a logic of places --- which is what
Chapter~\ref{ch:sci} needs when a specimen affords exactly one
measurement and hypotheses must therefore be about kinds.
And the comm rule is one definitional step from a crossover operator,
which is the observation Part~\ref{part:origins} builds on.

\section{Organization of this part}

Chapter~\ref{ch:gslt} gives the definition and argues for its two
halves --- why lambda theories rather than Lawvere theories, why the
rewrites are adjoined rather than enriched.
Chapters~\ref{ch:igslt} and~\ref{ch:cigslt} give the two strengthenings,
with the examples that separate them.
Chapter~\ref{ch:catgslt} assembles the category, whose morphisms preserve
bisimulation over context-labelled transitions.
Chapters~\ref{ch:costmonad} and~\ref{ch:histmonad} give the two
constructions on it that the rest of the book uses constantly: the cost
monad, which meters a theory, and the history monad, which logs it.
Chapter~\ref{ch:carho} exhibits the cost construction concretely, as a
metered rho calculus in which token stacks are first-class terms and
acceptance is a linear proof rather than a run to completion.
Chapter~\ref{ch:vtok} asks the question that metering raises and does not
answer --- when do many kinds of resource admit a single number? --- and
finds that the condition is exactly the absence of arbitrage.
Chapters~\ref{sec:lts-hml} and~\ref{sec:oslf-inputs} turn to the logic:
the labelled transition system a theory induces, the context-decorated
modal logic over it, and an accounting of what still has to be supplied
before that logic can be said to have been \emph{manufactured} rather
than postulated.
Chapter~\ref{ch:hypercube} closes the part with the construction that
does the manufacturing for type systems: OSLF and the Hypercube, which
read a theory's own rewrite rules and return a family of type systems for
it.

% ===================================================================
